% Source-derived chapter 06: Formula for divisor classes
% Original source: rigid-local-systems-multiplicative-eigenvalue-problem
\chapter{Formula for divisor classes}
\label{rigid-local-systems-multiplicative-eigenvalue-problem-chapter-06}
\source{rigid-local-systems-multiplicative-eigenvalue-problem}

\begin{remark}[Source section]
\label{rigid-local-systems-multiplicative-eigenvalue-problem-chapter-06-source}
\source{rigid-local-systems-multiplicative-eigenvalue-problem}
\source{rigid-local-systems-multiplicative-eigenvalue-problem}
This chapter reproduces the corresponding section of the retrieved
LaTeX source, including its definitions, statements, examples, and
proofs. It has not yet been translated into Lean.
\notready
\end{remark}

% The source section title was: Formula for divisor classes
\label{enume}
\source{rigid-local-systems-multiplicative-eigenvalue-problem}
Let $\vec{J}=(J^1,\dots,J^s)$ be an $s$-tuple of subsets  of $[n]$ each of cardinality $r$. Recall from Definition \ref{cyclist} that the effective cycle $E=C(d,r,\mn,n,\vec{J})$ on $\Par_{n,\mn,S}$ is defined as a pushforward of
$\Omega(d,r,\mn,n,\vec{J})$ (see Definition  \ref{mindful}).  We will compute $E$ when the codimension of this class is one, i.e.,
\begin{equation}\label{cdc}
\source{rigid-local-systems-multiplicative-eigenvalue-problem}
dn +\deg(\mn)r +r(n-r)-\sum_{i=1}^s |\sigma_{J^i}|=-1.
\end{equation}
We will use the following enumerative computation to prove Theorem \ref{brahms} in Section \ref{rumble}. It is also used to compare KZ local systems to strange duals of F-line bundles in Section \ref{EqKZ}. Recall the generalization of Gromov-Witten numbers given in  Definition \ref{GWgen}.
\begin{proposition}
\source{rigid-local-systems-multiplicative-eigenvalue-problem}
\notready\label{enumerative}
\source{rigid-local-systems-multiplicative-eigenvalue-problem}
Assume the codimension condition \eqref{cdc} and
write
$$\mathcal{O}(E)=\mathcal{B}(\vec{\lambda},\ell).$$
Write $\lambda^i=\sum_{b=1}^{n-1} c_{i}^{b}\omega_b$, where $\omega_b$ is the $b$th fundamental weight.
\begin{enumerate}
\item[(A)] Let $J^{s+1}=\{1,n-r+1,n-r+2,\dots,n-1\}$. Then
$$\ell=\langle\sigma_{J^1}, \sigma_{J^2},\dots,\sigma_{J^s},\sigma_{J^{s+1}}\rangle_{d+1,D}.$$
\item[(B)] $c_{i}^{b}=0$ if $b\not\in J^i$, or $b+1\in J^i$. If $b\in J^i$ and $b+1\not\in J^i$, define subsets $J'^1,\dots,J'^s$ of $[n]$
 each of cardinality $r$ as follows: $J'^k=J^k$ if $k\neq i$, and $J'^i=(J^i-\{b\})\cup \{b+1\}$. Then
$$c_i^{b}= \langle\sigma_{J'^1}, \sigma_{J'^2},\dots,\sigma_{J'^s}\rangle_{d,D}.$$
\item[(C)] If $1\leq b<n-1$ and $i\in[s]$, then $c^b_i\neq 0$ implies $c^{b+1}_i=0$.
\item[(D)] If $c^1_i\neq 0$ then $\ell=\sum_{b=1}^{n-1} c^b_i$.  Similarly if $c^{n-1}_i\neq 0$ then $\ell=\sum_{b=1}^{n-1} c^b_i$.
\end{enumerate}
\end{proposition}
We prove this proposition in the rest of this section. The  proof is modelled after the proofs of analogous enumerative counts in \cite{BHermit,BKiers,Kiers}.
\subsection{Formula (B) implies Formula (A)}
Add a new point $p_{s+1}$ and let $J'^{s+1}=\{n-r+1,\dots, n\}$, and $J'^i=J^i$ for $1\leq i \leq s$. Then, by Lemma \ref{propagation} below,
$$\mathcal{O}(C(d,r,\mathcal{O}(D),n,\vec{J'}))=\mathcal{D}^{\ell}\tensor\tensor_{i=1}^s\ml_{\lambda^i,p_i}\tensor \ml_{\vec{0},p_{s+1}}.$$
Now, by Proposition \ref{easily},
$$\operatorname{ISh}_{p^{s+1}}^*\mathcal{O}(C(d,r,\mathcal{O}(D),n,\vec{J'}))=\mathcal{D}^{\ell}\tensor\tensor_{i=1}^s\ml_{\lambda^i,p_i}\tensor\ml_{\lambda{s+1},p_{s+1}}$$
where $\lambda^{s+1}=(0,0,\dots,-\ell)$.

 $\operatorname{Sh}_{p^{s+1}}$ maps $C(d,r,\mathcal{O}(D),n,\vec{J'})$ isomorphically to $C(d,r,\mathcal{O}(D-1),n,\vec{K})$,
where $K^i=J^i$ for $1\leq i \leq s$ and $K^{s+1}=\{n-r,n-r+1,\dots,n-1\}$. Therefore,
$$\operatorname{ISh}_{p^{s+1}}^*C(d,r,\mathcal{O}(D),n,\vec{J'})=C(d,r,\mathcal{O}(D-1),n,\vec{K}).$$

Therefore we obtain the formula
$$\mathcal{O}(C(d,r,\mathcal{O}(D-1),n,\vec{K}))= \mathcal{D}^{\ell}\tensor\tensor_{i=1}^s\ml_{\lambda^i,p_i}\tensor\ml_{\lambda^{s+1},p_{s+1}}$$
where $\lambda^{s+1}=(0,0,\dots,-\ell)$. We have another formula for $\ell= \lambda_{n-1}^{s+1}-\lambda_{n}^{s+1}$ using  (B):
\begin{equation}\label{ABL1}
\source{rigid-local-systems-multiplicative-eigenvalue-problem}
\ell= \lambda_{n-1}^{s+1}-\lambda_{n}^{s+1}=\langle\sigma_{I^1}, \sigma_{I^2},\dots,\sigma_{I^s},\sigma_{I^{s+1}}\rangle_{d,D-1}
\end{equation}
with $I^j=K^j$ for $1\leq j\leq s$ and $I^{s+1}=\{n-r,\dots,n-2,n\}$. We can use the shift equalities, see Remark \ref{stuffed},  to get the stated formula for $\ell$.

\begin{remark}
\source{rigid-local-systems-multiplicative-eigenvalue-problem}
\notready\label{earlyAM}
\source{rigid-local-systems-multiplicative-eigenvalue-problem}
We also get formulas for $\ell$ without adding new points.  Pick any $p_i$. then $\ell-(\lambda^i_{1}-\lambda^i_{r})=0$ if either $1\in J^i$ or $n\not\in J^i$. If $1\not\in J^i$ and $n\in J^i$, then
\begin{equation}\label{ABL4}
\source{rigid-local-systems-multiplicative-eigenvalue-problem}
\ell-\sum_{b=1}^{n-1}c^b_i=\ell-(\lambda^i_{1}-\lambda^i_{r})=\langle\sigma_{J'^1}, \sigma_{J'^2},\dots,\sigma_{J'^s}\rangle_{d+1,D}
\end{equation}
where $J'^j=J^j$ for $j\neq i$ and $J'_i= (J^{i}-\{n\})\cup \{1\}$. Parallel to \eqref{ABL1}, we can shift this to
$$\ell-(\lambda^i_{1}-\lambda^i_{r})=\langle\sigma_{I^1}, \sigma_{I^2},\dots,\sigma_{I^s}\rangle_{d,D-1}$$
where $I^j=J^j$ for $j\neq i$ and $I^i= \{a\mid a=n\text{ or } a+1\in J^i\}$. Recall that $\lambda^i_1-\lambda^i_r=\sum_{b=1}^{n-1} c^b_i$.
\end{remark}

\begin{lemma}
\source{rigid-local-systems-multiplicative-eigenvalue-problem}
\notready\label{propagation}
\source{rigid-local-systems-multiplicative-eigenvalue-problem}
 Let $p_{s+1}\in \Bbb{P}^1-S$, and set $I^j=J^j$ for $\leq j\leq s$ and  $I^{s+1}=\{n-r+1,n-r+2,\dots,n\}$ so that $\sigma_{I^{s+1}}=1\in H^0(\Gr(r,n),\Bbb{Z})$. Let $S'=S\cup\{p_{s+1}\}$ and $\tau: \Par_{n,\mn,S'}\to \Par_{n,\mn,S}$ be the natural map which forgets the quasi-parabolic structure at $p_{s+1}$. Then,
 $$\tau^*C(d,r,\mn,n,\vec{J})=C(d,r,\mn,n,\vec{I}).$$
 \end{lemma}
\subsection{Proof of Formula (B)}
Consider a curve $C\leto{\sim} \Bbb{P}^1$ in $\Par_{n,\mn,S}$ consisting of points $(\mw,\me,\gamma)$, with $\mw$ (general and fixed) and $\gamma$ fixed, and all flags $\me^p$ fixed except when $p=p_i$. The flag $\me^{p_i}$ is the only variable, and $E^{p_i}_k$ is fixed except for $E^{p_i}_{b}$ which varies in a one dimensional family, with $E^{p_i}_{b-1}\subset E^{p_i}_b\subset E^{p_i}_{b+1}$. Here $E^p_0=0$. The number $c^i_b$ is the degree of $\mathcal{O}(E)$ on this curve. If $b\not\in J^i$ or $b+1\in J^i$  the rank condition at $b$ is not needed in the set of conditions defining the closed Schubert variety $\Omega_{J^i}(\me_p)$. It is now easy to see that, in this case, $E$ does not meet $C$ (for general choices of the fixed elements above) and the degree is zero, as desired.

Now suppose $b\in J^i$ and $b+1\not\in J^i$ and $J'^1,\dots,J'^s$ be as in the statement of (B). Let  $\Omega=\Omega(d,r,\mn,n,\vec{J})$ and $\Omega'=\Omega(d,r,\mn,n,\vec{J}')$. Then $\Omega\subseteq \Omega'$ and let $\pi:\Omega'\to \Par_{n,\mn,S}$ be the natural projection. Now $\pi$ is a generically finite map of degree $m=\langle\sigma_{J'^1}, \sigma_{J'^2},\dots,\sigma_{J'^s}\rangle_{d,D}$.  The fibers of $\Omega'$ over points of $C$ are finite sets of subbundles of $\mw$. These sets do not vary with the point on $C$. Therefore $\pi^{-1}(C)$ is $m$ copies of $\Bbb{P}^1$. The variation in these $\Bbb{P}^1$'s is only in the  $E^{p_i}_b$ elements of the flag (at $p_i$). Each of these $\Bbb{P}^1$'s meet $\Omega\subseteq \Omega'$ in exactly one point (and transversally so) by \cite[Lemma 8.8]{BHermit}. Therefore, $c^i_b =\deg(\pi_*[\Omega]\cap C)= \deg([\Omega]\cap \pi^{-1}(C))=m$, as desired.
\subsection{Proof of parts (C) and (D)}
Part (C) is immediate from the formula in (B) because $c_i^b\neq 0$ implies $b+1\notin J^i$ hence
$c_i^{b+1}=0$. Part (D) follows similarly from Remark \ref{earlyAM}.

\begin{remark}
\source{rigid-local-systems-multiplicative-eigenvalue-problem}
\notready\label{ajabaja}
\source{rigid-local-systems-multiplicative-eigenvalue-problem}
If $r=1$ in Proposition \ref{enumerative}, then $\ell=1$ as well since the only numbers that appear in the small quantum products of projective spaces are $0$ and $1$. In this case, if $D=0$ by normalization, it is easy to see that $\mathcal{O}(E)=\mathcal{B}(\vec{\lambda},1)$ with $\vec{\lambda}=(\omega_{a_1},\dots,\omega_{a_s})$ with $n|\sum a_i$.
\end{remark}
\subsection{Divisors of F-line bundles}
Let $\ml=\mathcal{B}(\vec{\lambda},\ell)$ be an F-line bundle  on $\Par_{n,\mathcal{O},S}$. Write it as $\mathcal{O}(E)$ for a reduced irreducible divisor  $E\subseteq \Par_{n,\mathcal{O},S}$.
Write $\lambda^i=\sum_{b=1}^{n-1} \gamma_{i}^{b}\omega_b$.

\begin{lemma}
\source{rigid-local-systems-multiplicative-eigenvalue-problem}
\notready\label{powercord}
\source{rigid-local-systems-multiplicative-eigenvalue-problem}
The following are equivalent :
\begin{enumerate}
\item $\mathcal{L}$ is of the form $\mathcal{O}(D(a,j))$ for some choice of $(a,j)$ in Theorem \ref{Adagio}.
\item $\gamma^j_b$ =1 for some choice of $j,b$, or $\ell-\sum_{b=1}^{n-1} \gamma^j_b=1$ for some $j\in[s]$.
\end{enumerate}
\end{lemma}
\begin{proof}
It is easy to see that (1) implies (2) if $a\neq 1$, since  $\gamma^j_{a-1}=1$. If $a=1$, then we use Remark \ref{earlyAM}.

Assume (2). By Lemma \ref{divisorE}. $E$ is equal to an effective cycle $C(d,r,\mathcal{O},n,\vec{J})$ (Definition  \ref{cyclist})  of codimension one i.e., $dn  +r(n-r)-\sum_{i=1}^s |\sigma_{J^i}|=-1$. Applying Proposition \ref{enumerative}, we find an enumerative formula for $\ml$ in terms of the data of the cycle $C(d,r,\mathcal{O},n,\vec{J})$. If $\gamma^j_b=1$ for some $j,b$ take $I^i=J^i$ for $i\neq j$ and $I^j=(J^j-\{b\})\cup\{b+1\}$ and apply Theorem \ref{Adagio} for the cycle data $(d,r,\mn,I^1,\dots,I^s)$, and consider the pair $(b+1,j)$   If $\ell-\sum_{b=1}^{n-1} \gamma^j_b=1$, we can take $I^j=(J^j-\{n\})\cup\{1\}$, and the pair $(1,j)$, and replace $d$ by $d+1$ and use Remark \ref{earlyAM}.
\end{proof}
